\section{Literature comparison, checked on 21 November 2026}
\label{sec:literature}

The comparison distinguishes the row tensors $D_n,P_n$ in
$(\mathbb C^n)^{\otimes n}$ from the polynomial coefficient tensors
$S_{\det,n},S_{\mathrm{per},n}$ in
$(\mathbb C^{n^2})^{\otimes n}$. It also distinguishes CP border
rank, symmetric border rank, Waring rank, and cactus rank.
For a symmetric tensor, CP border rank is at most symmetric
border rank, which is at most Waring rank. A lower bound on
Waring rank alone does not lower-bound border rank; neither does
a lower bound on cactus rank alone without further hypotheses.

\subsection{The numerical comparisons actually verified}

Landsberg--Teitler \cite[Corollary~1.4, p.~2; \S9, pp.~13--14,
Table~3]{LT2010} give Waring-rank lower bounds $43$ for $\det_4$
and $40$ for $\mathrm{per}_4$. Their table gives border lower
bound $36$ for both. The distinction follows from their rank
definitions and the surrounding argument, even where PDF text
extraction loses the underline. The middle catalecticant has
rank $\binom{n}{\lfloor n/2\rfloor}^2$.

Farnsworth \cite[Theorem~1.6, p.~2]{Farnsworth2016} proves
symmetric border rank at least $38$ for $\det_4$.
Theorem~1.3 gives his general-$n$ determinant estimate for $n\geq5$;
Theorem~1.8 gives symmetric border lower bound $14$ for both
degree-three polynomials. His bound $38$ is not silently relabeled
as a CP lower bound. The CP bounds proved here imply symmetric
bounds and therefore improve the specific symmetric comparison
when their value exceeds $38$.

Conner--Harper--Landsberg
\cite[Theorem~1.2, \S1.2]{CHL2023} prove
\[
 \underline{\operatorname R}_{\mathrm{CP}}(S_{\det,3})=17.
\]
Their tensor is explicitly in $(\mathbb C^9)^{\otimes3}$.
Conner--Huang--Landsberg
\cite[Theorem~1.1 and the following identification, p.~3;
proof in \S5]{CHHL2023} prove
\[
 \underline{\operatorname R}_{\mathrm{CP}}(S_{\mathrm{per},3})=16.
\]
These are polynomial tensors, whereas the later small numbers
for size four in the following paragraph concern row tensors.

Han--Ju--Kim \cite[Theorem~1.1 and \S3]{HJK2026} prove
$\underline{\operatorname R}(D_4)=12$ over $\mathbb C$.
The inspected version is arXiv:2510.11051v2, revised on
18 August 2026; its version history records a repaired gap.
Their other paper \cite[Theorems~5.6--5.7]{HJK2025} gives
$\underline{\operatorname R}(P_4)=\operatorname R(P_4)=8$
in characteristic different from two. Neither $12$ nor $8$
is a bound for a polynomial tensor in sixteen variables.

Two further 2026 papers make this ambient distinction
especially important. Han--Song
\cite[Theorem~A, p.~3; Proposition~4.1, p.~10]{HS2026}
prove ordinary CP rank $\operatorname R(P_5)=16$ in
characteristic zero. Their Proposition~4.1 also states
$\underline{\operatorname R}(P_n)\geq
\binom{n}{\lfloor n/2\rfloor}+1$ over $\mathbb C$
and its subfields for $n\geq3$, using a limiting-spans
argument. Thus this row-family ``plus one'' phenomenon
has an explicit current source. Han
\cite[Theorem~A, p.~2]{Han2026} proves symmetric border
rank $2^{n-1}$ for the squarefree monomial
$x_1\cdots x_n$, equivalently for the symmetric row
tensor $P_n$, in characteristic zero. This theorem
concerns symmetric border rank; it does not identify
CP border rank for every $n$, and it does not concern
the permanent polynomial in $n^2$ variables.

The Waring and cactus improvements reviewed by
Boij--Teitler \cite[background, pp.~5--7; \S4, p.~15]{BT2020},
including the Waring lower bound $50$ for $\det_4$, do not
replace the border-rank comparison. Houston--Goucher--Johnston
\cite{HGJ2024} were also checked for their determinant formulas
and tensor/Waring-rank bounds; no competing polynomial border
lower bound at the target size was located in that text.

\subsection{Which part is standard substitution, and which part is evaluated here}

Landsberg--Micha\l ek \cite[Proposition~2.3, p.~3]{LM2018}
already allow every retained dimension $a'\leq a$ in
\[
 \underline{\operatorname R}(T)
 \geq a-a'+
 \min_{A'\in\operatorname{Gr}(a',A^*)}
       \underline{\operatorname R}(T|_{A'}).
\]
Thus passing to $j>1$ is not claimed to invent a substitution
principle. The contribution to verify is the evaluated
Grassmannian minimum for the present residual image, with
the lower certificates and the explicit degenerations giving
upper witnesses. The short proper-incidence proof is included
for the all-characteristics and higher-order formulation
required here.

Torus and Borel symmetry reductions are likewise established
tools. Findings P-01 and P-07 give direct coordinate proofs
with their characteristic hypotheses explicit; the exact
residual minima and certificates are the deductions of this run.

\subsection{The border-apolarity comparison}

Buczy\'nska--Buczy\'nski
\cite[Theorem~3.15, p.~15; Theorem~4.3, p.~18;
Corollary~4.5, p.~19]{BB2021} formulate border apolarity through
the multigraded Hilbert-scheme component
$\operatorname{Slip}_r(X)$ of limits of ideals of points.
The fixed-ideal theorem allows a connected solvable stabilizer
to fix a witnessing ideal. Conner--Harper--Landsberg
\cite[Theorems~2.3--2.4, \S2.3]{CHL2023} use fixed limiting
planes and fixed ideals to make this framework effective.

The current simultaneous version is
Buczy\'nska--Buczy\'nski
\cite[\S4.3, Theorem~4.7; field convention at the start of
\S1]{BB2026}. Over every algebraically closed field, for a
smooth projective toric $X$, a very ample line bundle $L$,
and a subspace $E$ of the corresponding section-dual space
with $\dim E=i$, it gives
\[
 [E]\in\sigma_{r,i}(X)
 \quad\Longleftrightarrow\quad
 \exists[I]\in\operatorname{Slip}_r(X),\quad I\subseteq\operatorname{Ann}(E).
\]
If a connected solvable group preserves $E$, one may choose
$I$ invariant under that group. This provides the appropriate
current framework for simultaneous residual decompositions.
It does not itself supply the numerical minima evaluated here.

Two hypotheses matter. Correct Hilbert function and apolar
containment alone do not replace membership in
$\operatorname{Slip}_r(X)$. Also, the torus in P-01 preserves
the ambient residual space $U$ but generally moves its chosen
subspace $W$. Consequently P-01 is proved by an explicit
degeneration and closedness, not by applying a fixed-ideal
theorem with a group that does not stabilize $W$.

The later Ma\'ndziuk paper
\cite[Theorems~1.1--1.3, p.~2; Theorems~3.5, 4.4, 5.2;
\S6]{Mandziuk2026} adds tests beyond weak border apolarity.
Its examples are a Sch\"onhage sum, a Leitner tensor product,
two skew-Coppersmith--Winograd copies, and a reconsideration
of $3\times3$ matrix multiplication. Dole\v z\'alek--Micha\l ek
\cite[Theorem~1.1; \S5.1; \S6.3]{DM2026} give nonlinear
Kronecker--Koszul equations; their \S6.3 leaves extensive
special-tensor tests, including determinant tensors, for
future work. These current papers were inspected; no
polynomial $4\times4$ border lower bound competing with
the present target was located in their stated results.

\paragraph{Finding FL-06 (a displayed survey bound is false as written).}
Dutta--Lysikov \cite[\S4.3, p.~29, immediately after Corollary~56]{DL2025}
display the universal lower bound $\binom{2n}{n}$ for border
Waring rank of both $\det_n$ and $\mathrm{per}_n$. Taken
literally, this is false at $n=2$.

\begin{proof}
For either sign, $ad\pm bc$ has a nonsingular symmetric
coefficient matrix of size four. The first catalecticant
therefore gives border Waring rank at least four. Conversely,
\[
 ad\pm bc=
 \frac{(a+d)^2-(a-d)^2\pm(b+c)^2\mp(b-c)^2}{4}.
\]
Over $\mathbb C$ the scalar coefficients can be absorbed
into the squared linear forms. Thus its Waring rank and
border Waring rank are four, whereas $\binom42=6$.
\end{proof}

The correct standard comparison verified above is the
middle-catalecticant bound
$\binom{n}{\lfloor n/2\rfloor}^2$. We do not infer which
correction the survey authors intended.

\paragraph{Scope of the search.}
Primary source statements, their ambient definitions, and
relevant citation chains were checked on 21 November 2026.
Queries included determinant/permanent polynomial tensors,
CP and symmetric border rank, the numbers $39$, $40$ and $42$,
Farnsworth, and border apolarity, including 2025--2026 work.
No published polynomial border lower bound of at least $39$
for either size-four polynomial was located in the sources
checked. This is a scoped literature-audit outcome, not a
proof of exhaustive priority. The numerical improvements
are asserted against the specific verified comparisons,
without promoting absence from a search into a historical
novelty theorem.
